\documentclass[10pt]{article}
\usepackage[margin=0.7in]{geometry}
\usepackage{booktabs}
\usepackage{longtable}
\usepackage{pdflscape}
\usepackage{CJKutf8}
\usepackage{array}
\usepackage{graphicx}
\usepackage{caption}
\usepackage{microtype}
\usepackage{amsmath}
\setlength{\parindent}{0pt}
\setlength{\parskip}{4pt}
\newcommand{\NA}{--}
\newcolumntype{L}[1]{>{\raggedright\arraybackslash}p{#1}}
\begin{document}
\begin{center}
{\Large\bfseries Supplementary Material}\\[4pt]
{\large Pathology-Knowledge-Guided Structured Annotation for Clinical Speech Audio}
\end{center}

This supplementary material provides additional analyses, implementation details, and reproducibility information for the main paper. It includes population- and label-level support statistics, validation and sensitivity analyses for the weak-label pipeline, uncertainty analyses, serialization-error analysis, model configurations, complete Qwen3-Omni prompts and demonstration sets, audio preprocessing details, and JSON parsing and repair procedures.

\section{Population-wise support of the 22-label schema}
The 150 expert-reference records contain 50 recordings from each clinical population (30 training, 10 validation, and 10 test records per population). Table~\ref{tab:pop_label_support} reports positive record-level support for each label. Table~\ref{tab:pop_dim_support} summarizes dimension coverage using both positive label instances and record-level presence.

\begin{longtable}{llrrrr}
\caption{Expert-reference label support across the three clinical populations. Each population contains 50 annotated records.}\label{tab:pop_label_support}\\
\toprule
Dimension & Label & Aphasia & Depression & Dysarthria & Total \\
\midrule
\endfirsthead
\toprule
Dimension & Label & Aphasia & Depression & Dysarthria & Total \\
\midrule
\endhead

Respiration & Onset difficulty & 24 & 0 & 3 & 27 \\
Respiration & Phonation interruption & 23 & 3 & 3 & 29 \\
Respiration & Audible inhalation & 1 & 0 & 1 & 2 \\
Phonation & Hoarse/rough voice & 12 & 4 & 33 & 49 \\
Phonation & Breathiness & 4 & 1 & 8 & 13 \\
Phonation & Vocal tremor & 0 & 1 & 11 & 12 \\
Phonation & Pressed voice & 0 & 6 & 8 & 14 \\
Phonation & Weak voice & 2 & 10 & 16 & 28 \\
Articulation & Imprecise consonants & 23 & 1 & 47 & 71 \\
Articulation & Reduced vowel fullness & 23 & 0 & 45 & 68 \\
Articulation & Hypernasality & 0 & 0 & 18 & 18 \\
Articulation & Hyponasality & 0 & 0 & 1 & 1 \\
Prosody & Abnormal stress & 6 & 0 & 16 & 22 \\
Prosody & Unnatural rhythm & 28 & 2 & 47 & 77 \\
Prosody & Word-by-word speech & 11 & 0 & 38 & 49 \\
Prosody & Fast speech & 6 & 0 & 0 & 6 \\
Prosody & Slow speech & 13 & 1 & 33 & 47 \\
Prosody & Monotone prosody & 0 & 0 & 22 & 22 \\
Prosody & Excessive fillers & 0 & 2 & 0 & 2 \\
Additional events & Sighing & 0 & 0 & 0 & 0 \\
Additional events & Crying-like voice & 0 & 1 & 0 & 1 \\
Additional events & Laughter & 0 & 1 & 0 & 1 \\
\bottomrule
\end{longtable}

\begin{table}[h]
\centering
\caption{Dimension-level support across clinical populations. Label-instance support sums positive labels within a dimension; record presence counts records containing at least one positive label in that dimension.}
\label{tab:pop_dim_support}
\scriptsize
\begin{tabular}{lrrrrrr}
\toprule
& \multicolumn{3}{c}{Label instances} & \multicolumn{3}{c}{Record presence} \\
\cmidrule(lr){2-4}\cmidrule(lr){5-7}
Dimension & Aphasia & Depression & Dysarthria & Aphasia & Depression & Dysarthria \\
\midrule
Respiration & 48 & 3 & 7 & 31 & 3 & 5 \\
Phonation & 18 & 22 & 76 & 13 & 17 & 36 \\
Articulation & 46 & 1 & 111 & 30 & 1 & 48 \\
Prosody & 64 & 5 & 156 & 31 & 4 & 48 \\
Additional events & 0 & 2 & 0 & 0 & 2 & 0 \\
\bottomrule
\end{tabular}
\end{table}

\section{Held-out validation of the frozen weak-label pipeline}
The frozen weak-label pipeline was applied without further adjustment to the 30 expert-annotated validation records. It achieved a Micro-Precision of 0.4609, Micro-Recall of 0.5048, and Micro-F1 of 0.4818. Tables~\ref{tab:weak_dim} and~\ref{tab:weak_label} provide dimension- and label-level agreement. Undefined quantities caused by zero denominators are shown as ``--''.

\begin{table}[h]
\centering
\caption{Dimension-wise agreement of frozen weak labels with expert references on the 30-record validation split. Support denotes positive reference label instances.}
\label{tab:weak_dim}
\scriptsize
\begin{tabular}{lrrrr}
\toprule
Dimension & Support & Precision & Recall & F1 \\
\midrule
Respiration & 8 & 0.1154 & 0.3750 & 0.1765 \\
Phonation & 21 & 0.1429 & 0.0476 & 0.0714 \\
Articulation & 28 & 0.6970 & 0.8214 & 0.7541 \\
Prosody & 48 & 0.5306 & 0.5417 & 0.5361 \\
Additional events & 0 & -- & -- & -- \\
\bottomrule
\end{tabular}
\end{table}

\begin{longtable}{llrrrrrr}
\caption{Label-wise agreement of frozen weak labels with expert references on the 30-record validation split.} \label{tab:weak_label}\\
\toprule
Dimension & Label & Ref. support & Pred. $+$ & Precision & Recall & F1 & $\kappa$ \\
\midrule
\endfirsthead
\toprule
Dimension & Label & Ref. support & Pred. $+$ & Precision & Recall & F1 & $\kappa$ \\
\midrule
\endhead

Respiration & Onset difficulty & 4 & 13 & 0.1538 & 0.5000 & 0.2353 & 0.0394 \\
Respiration & Phonation interruption & 3 & 13 & 0.0769 & 0.3333 & 0.1250 & -0.0448 \\
Respiration & Audible inhalation & 1 & 0 & -- & 0.0000 & 0.0000 & 0.0000 \\
Phonation & Hoarse/rough voice & 7 & 0 & -- & 0.0000 & 0.0000 & 0.0000 \\
Phonation & Breathiness & 3 & 0 & -- & 0.0000 & 0.0000 & 0.0000 \\
Phonation & Vocal tremor & 4 & 5 & 0.0000 & 0.0000 & 0.0000 & -0.1739 \\
Phonation & Pressed voice & 3 & 1 & 0.0000 & 0.0000 & 0.0000 & -0.0526 \\
Phonation & Weak voice & 4 & 1 & 1.0000 & 0.2500 & 0.4000 & 0.3662 \\
Articulation & Imprecise consonants & 14 & 13 & 0.7692 & 0.7143 & 0.7407 & 0.5291 \\
Articulation & Reduced vowel fullness & 11 & 13 & 0.7692 & 0.9091 & 0.8333 & 0.7235 \\
Articulation & Hypernasality & 3 & 7 & 0.4286 & 1.0000 & 0.6000 & 0.5349 \\
Articulation & Hyponasality & 0 & 0 & -- & -- & -- & -- \\
Prosody & Abnormal stress & 5 & 6 & 0.1667 & 0.2000 & 0.1818 & 0.0000 \\
Prosody & Unnatural rhythm & 14 & 12 & 0.6667 & 0.5714 & 0.6154 & 0.3243 \\
Prosody & Word-by-word speech & 11 & 8 & 0.7500 & 0.5455 & 0.6316 & 0.4670 \\
Prosody & Fast speech & 2 & 0 & -- & 0.0000 & 0.0000 & 0.0000 \\
Prosody & Slow speech & 12 & 14 & 0.5000 & 0.5833 & 0.5385 & 0.1892 \\
Prosody & Monotone prosody & 4 & 7 & 0.5714 & 1.0000 & 0.7273 & 0.6715 \\
Prosody & Excessive fillers & 0 & 2 & 0.0000 & -- & 0.0000 & 0.0000 \\
Additional events & Sighing & 0 & 0 & -- & -- & -- & -- \\
Additional events & Crying-like voice & 0 & 0 & -- & -- & -- & -- \\
Additional events & Laughter & 0 & 0 & -- & -- & -- & -- \\
\bottomrule
\end{longtable}

\section{Threshold sensitivity of the frozen weak-label rules}
A one-at-a-time sensitivity analysis perturbed each analyzable continuous threshold to $0.8\times$, $0.9\times$, $1.0\times$, $1.1\times$, and $1.2\times$ its frozen value while keeping all other rules fixed. The $1.0\times$ condition exactly reproduced the original 30-record validation predictions (Micro-F1 $=0.4818$). Table~\ref{tab:threshold_sensitivity} reports the resulting Micro-F1 ranges.

\begin{longtable}{p{0.30\textwidth}p{0.23\textwidth}rrp{0.17\textwidth}}
\caption{One-at-a-time threshold sensitivity on the 30-record validation split. Tested range denotes the minimum and maximum numerical threshold values across the $0.8\times$--$1.2\times$ perturbations.} \label{tab:threshold_sensitivity}\\
\toprule
Threshold & Affected label & Baseline & Tested range & Micro-F1 range \\
\midrule
\endfirsthead
\toprule
Threshold & Affected label & Baseline & Tested range & Micro-F1 range \\
\midrule
\endhead

\texttt{hard\_attack\_ratio\_gt} & Onset difficulty & 0.18 & 0.144--0.216 & 0.4589--0.4818 \\
\texttt{break\_frequency\_gt} & Phonation interruption & 0.2 & 0.16--0.24 & 0.4753--0.4840 \\
\texttt{break\_ratio\_gt} & Phonation interruption & 0.05 & 0.04--0.06 & 0.4818--0.4818 \\
\texttt{ast\_breath\_probability\_gt} & Audible inhalation & 0.15 & 0.12--0.18 & 0.4818--0.4818 \\
\texttt{hoarse\_cpps\_lt} & Hoarse/rough voice & 0.12 & 0.096--0.144 & 0.4775--0.4818 \\
\texttt{hoarse\_hnr\_lt} & Hoarse/rough voice & 4 & 3.2--4.8 & 0.4818--0.4818 \\
\texttt{hoarse\_shimmer\_ge} & Hoarse/rough voice & 1.1 & 0.88--1.32 & 0.4818--0.4818 \\
\texttt{breathy\_cpps\_lt} & Breathiness & 0.1 & 0.08--0.12 & 0.4818--0.4818 \\
\texttt{breathy\_hnr\_lt} & Breathiness & 6 & 4.8--7.2 & 0.4818--0.4818 \\
\texttt{tremor\_jitter\_ge} & Vocal tremor & 0.05 & 0.04--0.06 & 0.4732--0.4885 \\
\texttt{tremor\_shimmer\_ge} & Vocal tremor & 1.1 & 0.88--1.32 & 0.4818--0.4818 \\
\texttt{squeezed\_loudness\_ge} & Pressed voice & 0.5 & 0.4--0.6 & 0.4818--0.4818 \\
\texttt{squeezed\_cpps\_ge} & Pressed voice & 0.18 & 0.144--0.216 & 0.4818--0.4818 \\
\texttt{squeezed\_hnr\_ge} & Pressed voice & 9 & 7.2--10.8 & 0.4818--0.4840 \\
\texttt{weak\_loudness\_lt} & Weak voice & 0.45 & 0.36--0.54 & 0.4749--0.4818 \\
\texttt{weak\_cpps\_lt} & Weak voice & 0.12 & 0.096--0.144 & 0.4818--0.4818 \\
\texttt{weak\_hnr\_lt} & Weak voice & 5 & 4--6 & 0.4818--0.4818 \\
\texttt{fast\_rate\_gt} & Fast speech & 3.3 & 2.64--3.96 & 0.4615--0.4818 \\
\texttt{slow\_rate\_lt} & Slow speech & 2.55 & 2.04--3.06 & 0.4476--0.4828 \\
\texttt{syllabic\_rate\_lt} & Word-by-word speech & 2.8 & 2.24--3.36 & 0.4135--0.4955 \\
\texttt{syllabic\_interval\_cv\_lt} & Word-by-word speech & 0.45 & 0.36--0.54 & 0.4135--0.4867 \\
\texttt{unnatural\_interval\_std\_gt} & Unnatural rhythm & 290 & 232--348 & 0.4779--0.4862 \\
\texttt{monotone\_f0\_cv\_lt} & Monotone prosody & 0.13 & 0.104--0.156 & 0.4775--0.4818 \\
\texttt{monotone\_loudness\_cv\_lt} & Monotone prosody & 0.85 & 0.68--1.02 & 0.4537--0.4818 \\
\texttt{stress\_ratio\_gt} & Abnormal stress & 8 & 6.4--9.6 & 0.4629--0.4818 \\
\texttt{filler\_rate\_gt} & Excessive fillers & 25 & 20--30 & 0.4775--0.4818 \\
\texttt{filler\_ratio\_gt} & Excessive fillers & 0.135 & 0.108--0.162 & 0.4775--0.4818 \\
\texttt{panns\_laugh\_probability\_gt} & Laughter & 0.15 & 0.12--0.18 & 0.4818--0.4818 \\
\texttt{panns\_cry\_probability\_gt} & Crying-like voice & 0.15 & 0.12--0.18 & 0.4818--0.4818 \\
\texttt{panns\_sigh\_probability\_gt} & Sighing & 0.15 & 0.12--0.18 & 0.4818--0.4818 \\
\texttt{vowel\_centroid\_std\_lt} & Reduced vowel fullness & 900 & 720--1080 & 0.4507--0.4818 \\
\texttt{hypernasal\_nor\_lower\_gt} & Hypernasality & 1.5 & 1.2--1.8 & 0.4796--0.4818 \\
\texttt{hypernasal\_hfr\_mean\_lt} & Hypernasality & 0.05 & 0.04--0.06 & 0.4818--0.4818 \\
\texttt{hyponasal\_nor\_upper\_lt} & Hyponasality & 1 & 0.8--1.2 & 0.4818--0.4818 \\
\bottomrule
\end{longtable}

\section{Test-set label support}
Table~\ref{tab:test_label_support} reports positive expert-reference support for all 22 labels on the 30-record test split, separated by clinical population. Each population contributes 10 test records. The three additional-event labels have zero positive test support and therefore do not receive dimension-level precision, recall, or F1 in the main paper.

\begin{longtable}{llrrrr}
\caption{Expert-reference label support on the 30-record test split by clinical population.}\label{tab:test_label_support}\\
\toprule
Dimension & Label & Aphasia & Depression & Dysarthria & Total \\
\midrule
\endfirsthead
\toprule
Dimension & Label & Aphasia & Depression & Dysarthria & Total \\
\midrule
\endhead

Respiration & Onset difficulty & 6 & 0 & 1 & 7 \\
Respiration & Phonation interruption & 6 & 0 & 1 & 7 \\
Respiration & Audible inhalation & 0 & 0 & 0 & 0 \\
Phonation & Hoarse/rough voice & 1 & 1 & 6 & 8 \\
Phonation & Breathiness & 0 & 0 & 2 & 2 \\
Phonation & Vocal tremor & 0 & 0 & 3 & 3 \\
Phonation & Pressed voice & 0 & 0 & 3 & 3 \\
Phonation & Weak voice & 0 & 0 & 5 & 5 \\
Articulation & Imprecise consonants & 7 & 0 & 10 & 17 \\
Articulation & Reduced vowel fullness & 7 & 0 & 10 & 17 \\
Articulation & Hypernasality & 0 & 0 & 3 & 3 \\
Articulation & Hyponasality & 0 & 0 & 1 & 1 \\
Prosody & Abnormal stress & 1 & 0 & 3 & 4 \\
Prosody & Unnatural rhythm & 8 & 0 & 10 & 18 \\
Prosody & Word-by-word speech & 2 & 0 & 9 & 11 \\
Prosody & Fast speech & 2 & 0 & 0 & 2 \\
Prosody & Slow speech & 2 & 0 & 7 & 9 \\
Prosody & Monotone prosody & 0 & 0 & 7 & 7 \\
Prosody & Excessive fillers & 0 & 0 & 0 & 0 \\
Additional events & Sighing & 0 & 0 & 0 & 0 \\
Additional events & Crying-like voice & 0 & 0 & 0 & 0 \\
Additional events & Laughter & 0 & 0 & 0 & 0 \\
\bottomrule
\end{longtable}

\begin{table}[h]
\centering
\caption{Dimension-level support on the 30-record test split.}
\label{tab:test_dim_support}
\scriptsize
\begin{tabular}{lrr}
\toprule
Dimension & Label-instance support & Record-presence support \\
\midrule
Respiration & 14 & 9 \\
Phonation & 21 & 10 \\
Articulation & 38 & 18 \\
Prosody & 51 & 18 \\
Additional events & 0 & 0 \\
\bottomrule
\end{tabular}
\end{table}

\section{Complete weak-label generation rules}
Table~\ref{tab:weak_rules} lists the final deterministic rules used by the frozen weak-label pipeline. The thresholds shown are the final post-calibration values, including substitutions made by the frozen \texttt{threshold\_final} wrapper.

\paragraph{Feature and model configuration.}
The general acoustic pipeline loads 16-kHz audio. Voice features are computed after \texttt{librosa.effects.split} with \texttt{top\_db=25}, frame length 2048, and hop length 512; a retained-voice proportion below 0.35 sets the voice-quality flag to low quality. openSMILE uses eGeMAPSv02 LLD/Functionals. Articulation features use \texttt{trim(top\_db=20)} and the STFT bands specified in Table~\ref{tab:weak_rules}. PANNs uses Cnn14 at 32 kHz with window 1024, hop 320, 64 mel bins, 50--14,000 Hz frequency limits, 527 output classes, and checkpoint \texttt{Cnn14\_mAP=0.431.pth}. Silero VAD is loaded from the project's local \texttt{silero-vad} copy.

\paragraph{Implementation-specific identifiers.}
The frozen source records PANNs event outputs by indices 16, 22, and 26. AST breathing-event classes are selected dynamically from the model configuration using the keyword set shown in Table~\ref{tab:weak_rules}.

\begin{landscape}
\scriptsize
\begin{longtable}{L{0.75in}L{1.05in}L{1.80in}L{2.45in}L{2.45in}}
\caption{Complete deterministic weak-label rules used by the frozen pipeline. Inclusive inequalities follow the implementation.}\label{tab:weak_rules}\\
\toprule
Dimension & Label & Evidence / feature & Frozen decision rule & Aggregation / implementation \\
\midrule
\endfirsthead
\toprule
Dimension & Label & Evidence / feature & Frozen decision rule & Aggregation / implementation \\
\midrule
\endhead

Respiration & Onset difficulty & RMS energy, onset detection, onset-window gradient & For each onset, count a hard onset when the energy gradient over the following 10 frames exceeds 0.012; assign the label when hard-onset ratio $>0.18$. & RMS frame length 256, hop 128; only onset windows with at least two frames enter the ratio. \\
Respiration & Phonation interruption & Silero VAD segments, \texttt{pyin} voiced flag, gap duration & Retain 30--200 ms gaps with local voiced-context mean $>0.3$; assign when break frequency $>0.2$ or break-duration ratio $>0.05$. & Silero uses minimum silence 30 ms and minimum speech 50 ms. The \texttt{pyin} context is evaluated around the gap midpoint ($\pm5$ frames). \\
Respiration & Audible inhalation & AST sigmoid probabilities and \texttt{id2label} & Select AST classes whose names contain \texttt{breath}, \texttt{gasp}, \texttt{pant}, \texttt{wheeze}, or \texttt{sniff}; assign when the maximum selected probability $>0.15$. & Whole waveform is passed to AST; aggregation is the maximum over matching output classes. \\
\midrule
Phonation & Hoarse/rough voice & CPPS p80, HNR p20, shimmer p80 & CPPS $\le0.12$ and (HNR $\le4.0$ or shimmer $\ge1.10$), with voice-quality flag \texttt{ok}. & Features are aggregated per audio after voice-region cropping; CPPS is a custom cepstral series. \\
Phonation & Breathiness & CPPS p80, HNR p20 & CPPS $\le0.10$ and HNR $\le6.0$, with voice-quality flag \texttt{ok}. & Same voice-feature aggregation as above. \\
Phonation & Vocal tremor & Jitter p80, shimmer p80 & Jitter $\ge0.05$ and shimmer $\ge1.10$, with voice-quality flag \texttt{ok}. & Same voice-feature aggregation as above. \\
Phonation & Pressed voice & Loudness p20, CPPS p80, HNR p20 & Loudness $\ge0.5$, CPPS $\ge0.18$, and HNR $\ge9.0$, with voice-quality flag \texttt{ok}. & Same voice-feature aggregation as above. \\
Phonation & Weak voice & Loudness p20, CPPS p80, HNR p20 & Loudness $\le0.45$ and (CPPS $\le0.12$ or HNR $\le5.0$), with voice-quality flag \texttt{ok}. & Same voice-feature aggregation as above; 0.45 is the final wrapper value. \\
\midrule
Articulation & Imprecise consonants & ZCR p95, spectral-flux standard deviation & Assign when ZCR p95 $<0.35$ or spectral-flux SD $<1.0$. & Computed on the trimmed 16-kHz waveform; 0.35 is the final wrapper value. \\
Articulation & Reduced vowel fullness & Spectral-centroid standard deviation & Assign when centroid SD $<900$. & Standard deviation across spectral-centroid frames; 900 is the final wrapper value. \\
Articulation & Hypernasality & Active-frame nasal/oral ratio (NOR) p20 and high-frequency ratio mean & NOR p20 $>1.5$ and mean high-frequency ratio $<0.05$. & STFT power bands: 200--500 Hz nasal, 500--2500 Hz oral, and $>2500$ Hz high-frequency; active frames have RMS above p20. \\
Articulation & Hyponasality & Active-frame NOR p95 & NOR p95 $<1.0$. & Same STFT bands and active-frame selection; at least 10 valid active frames are required by the feature routine. \\
\midrule
Prosody & Abnormal stress & Loudness p80/p20 ratio & Assign when $\mathrm{p80}/(\mathrm{p20}+10^{-6})>8.0$. & Loudness percentiles are read from the openSMILE Functionals output. \\
Prosody & Unnatural rhythm & RMS-peak interval SD, syllabic rate, interval CV & Interval SD $>290$ ms, or (syllabic rate $<2.8$ and interval CV $<0.45$). & RMS frame length 512, hop 256; peak height $=\max(\mathrm{RMS}_{dB})-25$ dB, distance 10, prominence 3. \\
Prosody & Word-by-word speech & Syllabic rate, interval CV & Syllabic rate $<2.8$ and interval CV $<0.45$. & RMS peaks define event times; rate is peak count divided by first-to-last peak duration; CV is interval SD divided by mean interval. \\
Prosody & Fast speech & Syllabic rate & Syllabic rate $>3.3$. & Same RMS-peak rate; no rhythm labels are returned when fewer than two peaks are detected. \\
Prosody & Slow speech & Syllabic rate & Syllabic rate $<2.55$. & Same RMS-peak rate; 2.55 is the final wrapper value. \\
Prosody & Monotone prosody & F0 normalized SD, loudness normalized SD & F0 normalized SD $<0.13$ and loudness normalized SD $<0.85$. & Values are read from openSMILE Functionals; 0.13 and 0.85 are final wrapper values. \\
Prosody & Excessive fillers & SenseVoiceSmall decoded text, fillers/minute, filler ratio & Assign when fillers/minute $>25$ or filler ratio $>0.135$. & After removing special tokens, count the fixed Mandarin filler lexicon \begin{CJK}{UTF8}{gbsn}\{嗯, 啊, 呃, 哎, 那个, 这个, 就是, 然后, 就\}\end{CJK}; rate uses audio duration and ratio uses cleaned character count. \\
\midrule
Additional events & Sighing & PANNs Cnn14 \texttt{clipwise\_output[26]} & Assign when probability $>0.15$. & One full-audio Cnn14 forward pass; no additional temporal pooling. \\
Additional events & Crying-like voice & PANNs Cnn14 \texttt{clipwise\_output[22]} & Assign when probability $>0.15$. & One full-audio Cnn14 forward pass; no additional temporal pooling. \\
Additional events & Laughter & PANNs Cnn14 \texttt{clipwise\_output[16]} & Assign when probability $>0.15$. & One full-audio Cnn14 forward pass; no additional temporal pooling. \\
\bottomrule
\end{longtable}
\end{landscape}




\section{Bootstrap uncertainty analyses}

\subsection{S4 versus S0}
We quantified uncertainty in the difference between the S4 label-aware construction and the S0 human-only baseline using 10,000 paired record-level bootstrap resamples of the common 30-record test set. The same resampled record indices were used for both systems, and training seeds were matched across S0 and S4. Differences are defined as S4 minus S0.

\begin{table}[h]
\centering
\caption{Paired record-level bootstrap comparison between S4 and S0. Positive Hamming-loss differences indicate higher loss for S4.}
\label{tab:s0_s4_bootstrap}
\scriptsize
\begin{tabular}{lrrrrrr}
\toprule
Metric & Observed $\Delta$ & Bootstrap mean $\Delta$ & 95\% CI & $\Delta>0$ & $\Delta<0$ & $\Delta=0$ \\
\midrule
Micro-F1     & 0.0052 & 0.0046 & [-0.0039, 0.0150] & 0.9661 & 0.0125 & 0.0214 \\
Exact Match  & 0.0334 & 0.0333 & [-0.0031, 0.0780] & 0.9125 & 0.0875 & 0.0000 \\
Hamming Loss & 0.0035 & 0.0036 & [-0.0019, 0.0117] & 0.8813 & 0.0442 & 0.0745 \\
\bottomrule
\end{tabular}
\end{table}

\subsection{SCSA versus Qwen3-Omni}
Cross-system uncertainty was estimated with 10,000 common-record bootstrap resamples of the same 30-record test set. Within each resample, metrics were computed for all five SCSA training runs and all five Qwen demonstration sets and then averaged within system. Differences are defined as SCSA minus Qwen3-Omni.

\begin{table}[h]
\centering
\caption{Common-record bootstrap comparison between SCSA S4 and Qwen3-Omni Q2.}
\label{tab:scsa_qwen_bootstrap}
\scriptsize
\begin{tabular}{lrrrrrr}
\toprule
Metric & Observed $\Delta$ & Bootstrap mean $\Delta$ & 95\% CI & $\Delta>0$ & $\Delta<0$ & $\Delta=0$ \\
\midrule
Precision & 0.3058 & 0.3289 & [0.2047, 0.4599] & 1.0000 & 0.0000 & 0.0000 \\
Recall    & -0.2378 & -0.2637 & [-0.3953, -0.1276] & 0.0002 & 0.9998 & 0.0000 \\
Micro-F1  & 0.0393 & 0.0468 & [0.0176, 0.0760] & 0.9504 & 0.0493 & 0.0003 \\
\bottomrule
\end{tabular}
\end{table}


\section{Flat-22 label-level error analysis}
The SCSA serialization ablation was further examined at label level over the five matched training seeds. Differences below are defined as Flat-22 minus five-dimensional serialization. The largest recall gain occurred for weak voice ($\Delta R=+0.8000$), while monotone prosody showed the largest F1 gain ($\Delta F1=+0.7692$). The largest precision decrease occurred for slow speech ($\Delta P=-0.3571$). Weak voice also produced the largest increase in false positives ($+10$), whereas hoarse/rough voice showed the largest reduction in false negatives ($-6$).

\begin{table}[h]
\centering
\caption{Largest label-level changes from five-dimensional to Flat-22 serialization.}
\label{tab:flat22_label_changes}
\small
\begin{tabular}{lll}
\toprule
Change & Label & Difference \\
\midrule
Largest recall increase & Weak voice & $+0.8000$ \\
Largest F1 increase & Monotone prosody & $+0.7692$ \\
Largest precision decrease & Slow speech & $-0.3571$ \\
Largest FP increase & Weak voice & $+10$ \\
Largest FN decrease & Hoarse/rough voice & $-6$ \\
\bottomrule
\end{tabular}
\end{table}

Across the 150 record--seed cases, the five-dimensional and Flat-22 representations yielded 49 and 6 exclusive exact matches, respectively, with no case in which both serializations were exactly correct. The corresponding exact-status overlap is shown in Table~\ref{tab:flat22_exact_overlap}. The most frequent failure combinations were phonation+articulation+prosody and respiration+articulation+prosody for the five-dimensional representation (five cases each), whereas phonation-only errors were most frequent for Flat-22 (11 cases). Among cases correctly matched only by the five-dimensional representation, weak voice was the most frequent additional false-positive label under Flat-22; no omitted false-negative label was observed in these cases.

\begin{table}[h]
\centering
\caption{Exact-match overlap between the two SCSA serialization variants across 150 record--seed cases.}
\label{tab:flat22_exact_overlap}
\small
\begin{tabular}{lrr}
\toprule
 & Flat-22 correct & Flat-22 incorrect \\
\midrule
Five-dimensional correct   & 0 & 49 \\
Five-dimensional incorrect & 6 & 95 \\
\bottomrule
\end{tabular}
\end{table}


\section{SCSA architecture and optimization configuration}
Table~\ref{tab:scsa_config} summarizes the architecture and training configuration used for SCSA. 

\begin{table}[h]
\centering
\caption{SCSA architecture and optimization configuration.}
\label{tab:scsa_config}
\scriptsize
\begin{tabular}{p{0.26\textwidth}p{0.66\textwidth}}
\toprule
Component & Final configuration \\
\midrule
CNN encoder & \texttt{captioning.models.cnn\_encoder.Cnn14Encoder}; PANNs checkpoint \texttt{Cnn14\_mAP=0.431.pth}; 2048-dimensional output; encoder frozen with batch-normalization layers kept in evaluation mode. \\
Recurrent encoder & Three-layer bidirectional GRU; hidden size 256 per direction; 512-dimensional output; dropout 0.5. \\
Transformer decoder & Two layers; model/token-embedding dimension 256; four attention heads; feed-forward dimension 1024; dropout 0.2. \\
Positional encoding & Fixed sinusoidal positional encoding with table length 100. \\
Vocabulary and target length & 84-token vocabulary; special tokens include \texttt{<pad>}, \texttt{<start>}, \texttt{<end>}, and \texttt{<unk>}; maximum target length 96 tokens. \\
Initialization and training seeds & The five S4 runs use training seeds 1--5; the seed affects trainable-parameter initialization, data order, dropout, and scheduled sampling, while the pretrained CNN14 encoder and fixed training files remain unchanged. \\
Loss & Token-masked mean label-smoothed cross-entropy; label smoothing 0.05. \\
Optimizer & Adam; learning rate $2\times10^{-4}$; weight decay $10^{-6}$; $\beta=(0.9,0.999)$; $\epsilon=10^{-8}$. \\
Learning-rate schedule & Exponential decay applied per iteration; 2,500 total iterations; 500 warm-up iterations; final learning rate $5\times10^{-7}$. \\
Training & Batch size 32; 20 epochs; gradient-norm clipping at 1.0; SpecAugment disabled. \\
Scheduled sampling & Enabled; linearly changed from 1.0 to 0.7 over training. \\
Checkpoint selection & Checkpoint with the minimum validation loss. \\
Main S4 inference & Beam search with beam size 3; maximum generation length 96; no length penalty. \\
Serialization-ablation inference & Greedy decoding (beam size 1), used for the matched five-dimensional versus Flat-22 comparison and kept distinct from the beam-decoded main S4 result. \\
\bottomrule
\end{tabular}
\end{table}


\section{Qwen3-Omni prompts and demonstration formatting}

\begin{table}[h]
\centering
\caption{Common Qwen3-Omni inference configuration used by the Q0--Q3 implementation.}
\label{tab:qwen_common_config}
\scriptsize
\begin{tabular}{p{0.28\textwidth}p{0.64\textwidth}}
\toprule
Item & Configuration \\
\midrule
Model & \texttt{Qwen3-Omni-30B-A3B-Instruct}. \\
Processor & \texttt{Qwen3OmniMoeProcessor}; \texttt{apply\_chat\_template(..., add\_generation\_prompt=True, tokenize=False)}. \\
Sampling / decoding & Deterministic beam search with \texttt{do\_sample=False} and beam size 3. \\
Generation length & Maximum generation length 96; length penalty 1.0. \\
Sampling controls & Sampling is disabled; no stochastic sampling is used. \\
Audio input & Audio is inserted as an audio content item after the corresponding text instruction. With audio standardization enabled, waveforms are loaded as 16-kHz mono audio, scaled by the 99.5th percentile of absolute amplitude when nonzero, and clipped to $[-1,1]$. \\
\bottomrule
\end{tabular}
\end{table}

\subsection{Q0: label-only zero-shot}
Q0 contains no system message. The test user message consists of the following text followed by the test audio:

\begin{CJK}{UTF8}{gbsn}
\begin{quote}\small
请根据输入音频判断是否包含下列标签。只依据音频内容选择标签，不要输出解释。

候选标签：\\{}
["起音困难", "发音中断", "明显吸气声", "沙哑/粗糙/嘶哑声", "气声", "声音颤抖", "挤压", "声音无力", "辅音模糊", "元音不饱满", "鼻音过重", "鼻音不足", "重音异常", "节奏不自然", "逐字发音", "语速过快", "语速过慢", "韵律单调", "填充词过多（嗯，啊）", "叹气", "哭腔", "笑声"]

只允许输出合法 JSON，不要输出解释文字。格式必须为：\\
\{\\
\quad "abnormal": true,\\
\quad "labels": []\\
\}
\end{quote}
\end{CJK}

\subsection{Q1 and Q2: pathology-guided five-dimensional prompt}
Q1 and Q2 use the same test-time pathology-guided instruction. Q1 is zero-shot, whereas Q2 prepends three audio--annotation demonstrations.

\begin{CJK}{UTF8}{gbsn}
\begin{quote}\small
你是一个专业的病理语音结构化标注系统。请根据输入音频完成二阶段判定，并严格输出固定 JSON schema。

阶段1：判断音频是否存在明确、可重复感知的病理语音异常。病理语音异常指由呼吸控制、发声质量、构音清晰度、韵律节奏或额外声学事件引起的稳定可感知异常；轻微噪声、录音设备差异或不确定现象不应标为异常。\\
阶段2：仅当 abnormal 为 true 时，从下列五个维度中选择具体标签；如果 abnormal 为 false，所有维度数组必须为空。

标签库：\\
- 呼吸: ["起音困难", "发音中断", "明显吸气声"]\\
- 发声: ["沙哑/粗糙/嘶哑声", "气声", "声音颤抖", "挤压", "声音无力"]\\
- 构音: ["辅音模糊", "元音不饱满", "鼻音过重", "鼻音不足"]\\
- 韵律: ["重音异常", "节奏不自然", "逐字发音", "语速过快", "语速过慢", "韵律单调", "填充词过多（嗯，啊）"]\\
- 额外声学事件: ["叹气", "哭腔", "笑声"]

只允许输出合法 JSON，不要输出解释文字。格式必须为：\\
\{\\
\quad "abnormal": true,\\
\quad "呼吸": [],\\
\quad "发声": [],\\
\quad "构音": [],\\
\quad "韵律": [],\\
\quad "额外声学事件": []\\
\}
\end{quote}
\end{CJK}

For each Q2 demonstration, the user message contains the following text followed immediately by the demonstration audio:

\begin{CJK}{UTF8}{gbsn}
\begin{quote}\small
请按照固定 schema 标注这个音频，只输出 JSON。
\end{quote}
\end{CJK}

The corresponding assistant message is the normalized five-dimensional expert annotation serialized as JSON. After three demonstration user/assistant pairs, the final user message contains the complete Q1/Q2 prompt above followed by the test audio.

\subsection{Q3: pathology-guided Flat-22 prompt}
Q3 uses the same pathology guidance, five dimensions, 22-label inventory, abnormal rule, and evidence rule as Q1/Q2. Its final output instruction is replaced by:

\begin{CJK}{UTF8}{gbsn}
\begin{quote}\small
只允许输出合法 JSON，不要输出解释文字。格式必须为：\\
\{\\
\quad "abnormal": true,\\
\quad "labels": []\\
\}
\end{quote}
\end{CJK}

For each Q3 demonstration, the user message contains the following text followed by the same demonstration audio used by Q2 for that seed:

\begin{CJK}{UTF8}{gbsn}
\begin{quote}\small
请按照 22 个候选标签标注这个音频，只输出包含 labels 的 JSON。
\end{quote}
\end{CJK}

The assistant demonstration target is obtained by normalizing the five-dimensional annotation and concatenating the retained labels in the fixed dimension order into the single \texttt{labels} array. Thus Q2 and Q3 differ in target serialization and the serialization-specific demonstration instruction, while using the same pathology guidance and matched demonstration audio.

\subsection{Demonstration sets}
Each Q2/Q3 demonstration set contains one expert-annotated training record from each clinical population. Q2 and Q3 use the identical ordered demonstration manifest within each seed.

\section{Audio preprocessing details}
Corpus segmentation and denoising were performed with the openSMILE-based preprocessing workflow. The processed corpus used by the downstream systems is stored as 16-kHz mono, signed 16-bit PCM WAV audio.

\begin{table}[h]
\centering
\caption{Confirmed corpus and model-specific audio preprocessing operations.}
\label{tab:audio_preprocessing}
\scriptsize
\begin{tabular}{p{0.22\textwidth}p{0.70\textwidth}}
\toprule
Stage & Confirmed processing \\
\midrule
Stored corpus &
16-kHz mono, signed 16-bit PCM WAV. \\

SCSA &
Manifest WAVs are loaded in mono and conditionally resampled to 16 kHz. The final S4 configuration uses no fixed-duration padding or cropping. \\

Qwen3-Omni &
For the retained Q0--Q2 runs, audio is loaded as 16-kHz mono waveform, divided by the 99.5th percentile of absolute amplitude when nonzero, and clipped to $[-1,1]$. No additional trimming, padding, cropping, or denoising is applied in this standardization step. \\

Weak-label voice features &
Audio is loaded as 16-kHz mono waveform and unit-peak normalized. Voice-feature extraction uses activity intervals obtained with \texttt{librosa.effects.split} at \texttt{top\_db=25}, \texttt{frame\_length=2048}, and \texttt{hop\_length=512}. \\

AST &
Audio is resampled to 16 kHz when needed; multichannel input is averaged to mono before feature extraction. \\

PANNs &
Audio is loaded and resampled in memory to 32-kHz mono immediately before Cnn14 inference. The Cnn14 front end uses a 1024-sample window, 320-sample hop, 64 mel bins, and a 50--14,000 Hz frequency range. \\

Articulation analysis &
Audio is loaded at 16 kHz and trimmed with \texttt{librosa.effects.trim(top\_db=20)}; clips shorter than 0.5 s after trimming are not analyzed by this module. \\

Glottal / break analysis &
Glottal-attack analysis uses 16-kHz audio with unit-peak normalization. The Silero-based phonation-break detector reads audio at 16 kHz with minimum silence and speech durations of 30 ms and 50 ms, respectively. \\
\bottomrule
\end{tabular}
\end{table}


\section{JSON parsing, normalization, and scanning repair}
Both systems apply deterministic pre-parsing before JSON validation. The validity and repair metrics are defined at fixed stages of this pipeline so that routine schema normalization is not counted as repair.

\begin{table}[h]
\centering
\caption{Confirmed JSON parsing and schema-recovery procedures.}
\label{tab:json_repair}
\scriptsize
\begin{tabular}{p{0.18\textwidth}p{0.34\textwidth}p{0.40\textwidth}}
\toprule
System & Pre-parsing and normalization & Scanning-repair fallback \\
\midrule
SCSA &
Decoded text is stripped; \texttt{<unk>}, \texttt{<pad>}, \texttt{<start>}, and \texttt{<end>} are removed; whitespace is collapsed; and two fixed regular-expression corrections for known formatting errors are applied before \texttt{json.loads}. A parseable dictionary is canonicalized by inserting the five dimensions, converting permitted scalar labels to one-element lists, exact-filtering the allowed vocabulary, and removing duplicates. &
The fallback is entered when JSON parsing fails or the parsed root is not a dictionary. For each dimension, the procedure scans the malformed text for the quoted dimension field and its following list region, extracts quoted strings, retains only exact allowed labels for that dimension, and removes duplicates. No synonym expansion, fuzzy matching, or label inference is used. A complete five-dimension object is constructed directly; there is no second \texttt{json.loads} success test. \\

Qwen3-Omni &
The inference path first selects text between the first \texttt{\{} and last \texttt{\}} when that candidate is parseable; otherwise the raw decoded text is retained. The parser again selects the brace-delimited candidate when available and applies \texttt{json.loads}. Parseable dictionaries are normalized by unwrapping the expected object when needed, mapping Flat-22 labels to dimensions, converting permitted scalar fields to lists, exact-filtering the vocabulary, and removing duplicates. &
The fallback is entered on JSON parse failure or a non-dictionary parsed root. For each dimension, the procedure scans the corresponding list region and retains exact occurrences of allowed labels. If a dimension region cannot be captured, it searches the whole candidate text for exact allowed-label strings. No synonym or fuzzy matching is used. The recovered object is then passed through the common schema normalizer. \\
\bottomrule
\end{tabular}
\end{table}

\paragraph{Validity and repair metrics.}
Pre-parsed JSON Valid Rate is measured after the model-specific deterministic pre-parsing described above and before schema normalization or scanning repair. Schema Repair Rate is the fraction of all evaluated outputs that enter the scanning-repair fallback:
\[
\mathrm{Schema\ Repair\ Rate}
=
\frac{\#\{\text{outputs entering scanning repair}\}}{N}.
\]
\end{document}
